\documentclass[12pt]{article}
\usepackage{amsmath, amssymb, amsthm}
\usepackage{geometry}
\usepackage{booktabs}
\usepackage{hyperref}
\usepackage{xcolor}
\geometry{margin=2.5cm}

\theoremstyle{definition}
\newtheorem{definition}{Definition}
\newtheorem{proposition}{Proposition}
\newtheorem{theorem}{Theorem}
\newtheorem{corollary}{Corollary}
\newtheorem{remark}{Remark}

\title{\textbf{Covariance of Empirical Periodogram Ordinates}\\
\large Mathematical Summary (Revised: Brillinger Convention)}
\author{GEMS TCO Project}
\date{\today}

\begin{document}
\maketitle
\tableofcontents
\newpage

%───────────────────────────────────────────────────────────────
\section{Setup and Notation}
%───────────────────────────────────────────────────────────────

Let $Y(\mathbf{s})$, $\mathbf{s} = (s_1, s_2) \in \{0,\ldots,M-1\}\times\{0,\ldots,N-1\}$,
be a real-valued, mean-zero, second-order stationary spatial process on
an $M\times N$ grid.
The autocovariance function is
\[
  C(\mathbf{h}) = \mathrm{Cov}(Y(\mathbf{s}),\, Y(\mathbf{s}+\mathbf{h})),
\]
and the spectral density satisfies
\[
  f(\boldsymbol{\omega}) = \sum_{\mathbf{h}} C(\mathbf{h})\,e^{-i\boldsymbol{\omega}\cdot\mathbf{h}},
  \qquad f(\boldsymbol{\omega}) \geq 0,
\]
with inverse
\[
  C(\mathbf{h}) = \frac{1}{MN}\sum_{\mathbf{k}} f(\boldsymbol{\omega}_{\mathbf{k}})\,
    e^{i\boldsymbol{\omega}_{\mathbf{k}}\cdot\mathbf{h}},
  \quad
  \boldsymbol{\omega}_{\mathbf{k}} = 2\pi\!\left(\frac{k_1}{M},\frac{k_2}{N}\right).
\]

\paragraph{Observation model with masking.}
In practice (GEMS TCO), cells without valid retrievals are
zero-imputed.  Let $m(\mathbf{s}) \in \{0,1\}$ be the binary observation
mask.  The zero-imputed field is
\[
  X(\mathbf{s}) \;=\; m(\mathbf{s})\,Y(\mathbf{s}).
\]
All formulas below are first derived for the complete field $Y$
(mask $m\equiv 1$) and then extended to $X$ in
Section~\ref{sec:masking}.


%───────────────────────────────────────────────────────────────
\section{Taper, Effective Weight, and Tapered DFT}
%───────────────────────────────────────────────────────────────

\subsection{Spatial taper}

Let $h(\mathbf{s}) = h_1(s_1)h_2(s_2)$ be a separable, non-negative taper
(e.g.\ Hamming window):
\[
  h_{(i,j)}
  \;=\; \Bigl(0.54 - 0.46\cos\tfrac{2\pi i}{M}\Bigr)\,
         \Bigl(0.54 - 0.46\cos\tfrac{2\pi j}{N}\Bigr).
\]

\subsection{Effective taper (taper $\times$ mask)}

Let $\mathcal{O}^{(q)}$ be the set of observed cells at snapshot $q$.
The \emph{effective weight} at cell $\mathbf{s}$ for snapshot $q$ is
\[
  g^{(q)}_{\mathbf{s}} \;=\; h_{\mathbf{s}}\cdot\mathbf{1}[\mathbf{s}\in\mathcal{O}^{(q)}].
\]
This combines the taper and the observation mask into a single weight;
zero-imputed missing cells contribute zero weight.
The normalisation constant is
\[
  H_q \;=\; \sum_{\mathbf{s}}\bigl(g^{(q)}_{\mathbf{s}}\bigr)^2
       \;=\; \sum_{\mathbf{s}\in\mathcal{O}^{(q)}} h_{\mathbf{s}}^2.
\]
When all cells are observed ($\mathcal{O}^{(q)} = $ full grid):
$H_q = \|h\|^2 \approx 0.375\,MN$ (Hamming) or $= MN$ (rectangular).

\begin{remark}[Why $H_q$ instead of $MN$]
Using $H_q$ rather than $MN$ as the normaliser ensures that the
periodogram is approximately unbiased for $f(\boldsymbol{\omega})$ regardless
of how many cells are observed.
If $1/(MN)$ were used, missing data would attenuate the periodogram
by the observation fraction $p = |\mathcal{O}^{(q)}|/MN$ --- exactly
the $p^2$ attenuation analysed in the companion spectral distortion paper.
\end{remark}

\subsection{Tapered DFT (J-vector, code convention)}

The tapered DFT at physical frequency $\boldsymbol{\omega}_{\mathbf{k}}$ for snapshot $q$ is
\[
\boxed{
  J^{(q)}(\boldsymbol{\omega}_{\mathbf{k}})
  \;=\; \frac{1}{2\pi\sqrt{H_q}}\,
    \mathcal{F}_{2D}\!\bigl[g^{(q)}\odot X^{(q)}\bigr](\mathbf{k}),
}
\]
where $\mathcal{F}_{2D}[\cdot](\mathbf{k}) = \sum_{\mathbf{s}}(\cdot)_{\mathbf{s}}\,
e^{-i2\pi(k_1 s_1/M + k_2 s_2/N)}$ is the standard (non-unitary) 2-D DFT,
and the physical frequencies are
$\boldsymbol{\omega}_{\mathbf{k}} = (2\pi k_1/(M\Delta_1),\; 2\pi k_2/(N\Delta_2))$
for grid spacings $\Delta_1, \Delta_2$.
Expanding explicitly:
\[
  J^{(q)}(\boldsymbol{\omega}_{\mathbf{k}})
  \;=\; \frac{1}{2\pi\sqrt{H_q}}\sum_{\mathbf{s}\in\mathcal{O}^{(q)}}
    h_{\mathbf{s}}\,X^{(q)}(\mathbf{s})\,
    e^{-i\boldsymbol{\omega}_{\mathbf{k}}\cdot\mathbf{s}\,\Delta},
\]
where $\boldsymbol{\omega}_{\mathbf{k}}\cdot\mathbf{s}\,\Delta = \omega^{\mathrm{lat}}_{k_1} s_1\Delta_1
+ \omega^{\mathrm{lon}}_{k_2} s_2\Delta_2$.

\begin{remark}[Asymptotic distribution]
Asymptotically, by Brillinger (1975, Thm.~4.4.1):
\[
  J^{(q)}(\boldsymbol{\omega}_{\mathbf{k}}) \;\sim\; \mathcal{CN}\!\Bigl(0,\,\tfrac{f(\boldsymbol{\omega}_{\mathbf{k}})}{(2\pi)^2}\Bigr).
\]
Writing $J^{(q)} = a_{\mathbf{k}} + ib_{\mathbf{k}}$:
$a_{\mathbf{k}}, b_{\mathbf{k}} \sim \mathcal{N}(0,\, f(\boldsymbol{\omega}_{\mathbf{k}})/(2(2\pi)^2))$,
independently, so
$\mathbb{E}[I^{(q)}_{\mathbf{k}}] = \mathbb{E}[|J^{(q)}(\boldsymbol{\omega}_{\mathbf{k}})|^2]
\approx f(\boldsymbol{\omega}_{\mathbf{k}})/(2\pi)^2$.
The periodogram estimates $f/(2\pi)^2$, not $f$ directly;
the factor $(2\pi)^2$ is part of the spectral density convention used
in \texttt{expected\_periodogram\_fft\_tapered}.
\end{remark}

\subsection{Sample cross-periodogram}

For two snapshots $q, r$:
\[
  \hat{I}^{(q,r)}(\boldsymbol{\omega}_{\mathbf{k}})
  \;=\; J^{(q)}(\boldsymbol{\omega}_{\mathbf{k}})\,\overline{J^{(r)}(\boldsymbol{\omega}_{\mathbf{k}})}.
\]
The empirical expectation over $N_{\mathrm{iter}}$ realisations is
\[
  \bar{I}^{(q,r)}_{\mathrm{emp}}(\boldsymbol{\omega}_{\mathbf{k}})
  \;=\; \frac{1}{N_{\mathrm{iter}}}\sum_{\ell=1}^{N_{\mathrm{iter}}}
    \hat{I}^{(q,r)}_\ell(\boldsymbol{\omega}_{\mathbf{k}}).
\]

\subsection{DFT of the taper (auxiliary, for spectral-domain formulas)}

For the spectral-domain formula of $\tilde{C}_g$
(Section~\ref{sec:Ctilde_spectral}), define
\[
  W^{(q)}[\mathbf{m}]
  \;=\; \sum_{\mathbf{s}} g^{(q)}_{\mathbf{s}}\,e^{-i2\pi\mathbf{m}\cdot\mathbf{s}/(MN)}
  \;=\; \mathrm{FFT}(g^{(q)})[\mathbf{m}].
\]
Note that $W^{(q)} = \mathrm{FFT}(h \cdot m^{(q)})$ combines the
taper and mask spectra; when $m^{(q)}\equiv 1$, this reduces to
$\mathrm{FFT}(h)$.


%───────────────────────────────────────────────────────────────
\section{The Effective-Taper Cross-Covariance Kernel $\tilde{C}_g$}
\label{sec:Ctilde}
%───────────────────────────────────────────────────────────────

\subsection{Definition}

\begin{definition}[Effective-taper cross-covariance kernel]
\label{def:Ctilde}
For snapshot $q$ with effective weight $g^{(q)}$ and normalisation $H_q$,
and for physical frequencies $\boldsymbol{\omega}_{\mathbf{j}},\boldsymbol{\omega}_{\mathbf{k}}$, define
\[
\boxed{
  \tilde{C}_{g^{(q)}}(\boldsymbol{\omega}_{\mathbf{j}},\boldsymbol{\omega}_{\mathbf{k}})
  \;=\; \frac{1}{(2\pi)^2 H_q}
    \sum_{\mathbf{u},\mathbf{v}}
    g^{(q)}_{\mathbf{u}}\,g^{(q)}_{\mathbf{v}}\,C(\mathbf{u}-\mathbf{v})\,
    e^{-i(\boldsymbol{\omega}_{\mathbf{j}}\cdot\mathbf{u}\Delta\,-\,\boldsymbol{\omega}_{\mathbf{k}}\cdot\mathbf{v}\Delta)}.
}
\]
\end{definition}

This equals the expected cross-periodogram between the DFT coefficients:
\[
  \tilde{C}_{g^{(q)}}(\boldsymbol{\omega}_{\mathbf{j}},\boldsymbol{\omega}_{\mathbf{k}})
  \;=\; \mathbb{E}_Y\!\left[J^{(q)}(\boldsymbol{\omega}_{\mathbf{j}})\,
                             \overline{J^{(q)}(\boldsymbol{\omega}_{\mathbf{k}})}\right].
\]
The mask is already incorporated in $g^{(q)} = h\cdot m^{(q)}$,
so no separate masking treatment is needed at this level.
For notational convenience we drop the superscript $(q)$ and write
$\tilde{C}_g$ when the snapshot is clear from context.

\begin{remark}[Physical interpretation]
$\tilde{C}_g(\boldsymbol{\omega}_{\mathbf{j}},\boldsymbol{\omega}_{\mathbf{k}})$ measures spectral leakage
from true frequency $\boldsymbol{\omega}_{\mathbf{k}}$ into observed frequency
$\boldsymbol{\omega}_{\mathbf{j}}$ through the effective window $g$.
For a rectangular window with full observations ($g\equiv 1$,
$H_q = MN$): $\tilde{C}_g(\omega_{\mathbf{j}},\omega_{\mathbf{k}})
= f(\boldsymbol{\omega}_{\mathbf{j}})/(2\pi)^2\cdot\delta_{\mathbf{j}=\mathbf{k}}$.
\end{remark}

\subsection{Spectral-domain equivalent}
\label{sec:Ctilde_spectral}

Substituting $C(\mathbf{h}) = (MN)^{-1}\sum_{\mathbf{m}} f(\boldsymbol{\omega}_{\mathbf{m}})
e^{i\boldsymbol{\omega}_{\mathbf{m}}\cdot\mathbf{h}\Delta}$ into Definition~\ref{def:Ctilde}:

\begin{proposition}[Spectral convolution representation]
\label{prop:spectral_conv}
\[
  \tilde{C}_g(\boldsymbol{\omega}_{\mathbf{j}},\boldsymbol{\omega}_{\mathbf{k}})
  \;=\; \frac{1}{(2\pi)^2(MN)H_q}\sum_{\mathbf{m}} f(\boldsymbol{\omega}_{\mathbf{m}})\,
    W^{(q)}[\mathbf{j}-\mathbf{m}]\,\overline{W^{(q)}[\mathbf{k}-\mathbf{m}]}.
\]
\end{proposition}

\textbf{Covariance-domain form is preferred for code} (Definition~\ref{def:Ctilde}):
it uses $C(\mathbf{u}-\mathbf{v})$ and $g^{(q)}$ directly without computing $W^{(q)}$.
The spectral-domain form (Proposition~\ref{prop:spectral_conv}) is useful
for theoretical analysis and for connecting to Sykulski et al.\ (2019).

\subsection{Pseudo cross-spectral density}

For a real-valued process, $J(-\boldsymbol{\omega}) = \overline{J(\boldsymbol{\omega})}$, so
\[
  \tilde{P}_g(\boldsymbol{\omega}_{\mathbf{j}},\boldsymbol{\omega}_{\mathbf{k}})
  \;=\; \mathbb{E}_Y\!\left[J^{(q)}(\boldsymbol{\omega}_{\mathbf{j}})\,J^{(q)}(\boldsymbol{\omega}_{\mathbf{k}})\right]
  \;=\; \tilde{C}_g(\boldsymbol{\omega}_{\mathbf{j}},\,-\boldsymbol{\omega}_{\mathbf{k}}).
\]


%───────────────────────────────────────────────────────────────
\section{Exact Covariance of the Empirical Periodogram}
\label{sec:exact_cov}
%───────────────────────────────────────────────────────────────

\subsection{Gaussian case (Wick's theorem)}

\begin{theorem}[Periodogram covariance via $\tilde{C}_g$]
\label{thm:cov_gaussian}
If $Y$ is Gaussian and the effective weight is $g^{(q)}$, then for any $\mathbf{j},\mathbf{k}$:
\[
\boxed{
  \mathrm{Cov}(I^{(q)}_{\mathbf{j}},\, I^{(q)}_{\mathbf{k}})
  \;=\;
  \bigl|\tilde{C}_g(\boldsymbol{\omega}_{\mathbf{j}},\boldsymbol{\omega}_{\mathbf{k}})\bigr|^2
  \;+\;
  \bigl|\tilde{C}_g(\boldsymbol{\omega}_{\mathbf{j}},\,-\boldsymbol{\omega}_{\mathbf{k}})\bigr|^2.
}
\]
\end{theorem}

\begin{proof}
By Wick's theorem (Isserlis):
$\mathrm{Cov}(|J^{(q)}_{\mathbf{j}}|^2, |J^{(q)}_{\mathbf{k}}|^2)
= |\mathbb{E}[J^{(q)}_{\mathbf{j}}\overline{J^{(q)}_{\mathbf{k}}}]|^2
+ |\mathbb{E}[J^{(q)}_{\mathbf{j}}J^{(q)}_{\mathbf{k}}]|^2
= |\tilde{C}_g|^2 + |\tilde{P}_g|^2$. \qed
\end{proof}

\begin{corollary}[Simplified form for positive frequencies]
\label{cor:pos_freq}
For a real-valued process and positive Fourier frequencies
$\mathbf{j},\mathbf{k}\neq\mathbf{0}$ with $\mathbf{j}\neq-\mathbf{k}$:
\[
\boxed{
  \mathrm{Cov}(I^{(q)}_{\mathbf{j}},\, I^{(q)}_{\mathbf{k}})
  \;\approx\;
  2\,\bigl|\tilde{C}_g(\boldsymbol{\omega}_{\mathbf{j}},\boldsymbol{\omega}_{\mathbf{k}})\bigr|^2.
}
\]
\end{corollary}

This is the form used in the code (Brillinger, 1975, \S7.6).

\subsection{Non-Gaussian case}

\[
  \mathrm{Cov}(I^{(q)}_{\mathbf{j}},\, I^{(q)}_{\mathbf{k}})
  \;=\;
  \bigl|\tilde{C}_g(\boldsymbol{\omega}_{\mathbf{j}},\boldsymbol{\omega}_{\mathbf{k}})\bigr|^2
  +
  \bigl|\tilde{C}_g(\boldsymbol{\omega}_{\mathbf{j}},\,-\boldsymbol{\omega}_{\mathbf{k}})\bigr|^2
  + \kappa_{\mathbf{j},\mathbf{k}},
\]
where $\kappa_{\mathbf{j},\mathbf{k}}$ is the trispectrum contribution
(Section~\ref{sec:nongauss}).

\subsection{Hamming taper: spectral window width}

For the Hamming window, $W^{(q)}[\mathbf{m}]$ (when $m^{(q)}\equiv 1$) has
support only at $\mathbf{m}\in\{\mathbf{0},\pm\mathbf{e}_1,\pm\mathbf{e}_2\}$
(3 non-negligible values per axis), so $\tilde{C}_g$ mixes power over
roughly $\pm 1$ frequency bin.
For the covariance estimators, we evaluate Definition~\ref{def:Ctilde}
directly (covariance domain), avoiding explicit computation of $W^{(q)}$.


%───────────────────────────────────────────────────────────────
\section{Masking Effect: $H_q$ Normalisation vs.\ $MN$}
\label{sec:masking}
%───────────────────────────────────────────────────────────────

Because the mask $m^{(q)}$ is already baked into the effective weight
$g^{(q)} = h\cdot m^{(q)}$, there is no separate masking step in the
DFT or in the $\tilde{C}_g$ formula.
The masking enters solely through:
\begin{enumerate}
  \item the sum in $\tilde{C}_g$ running only over observed cells
        ($g^{(q)}_{\mathbf{s}} = 0$ for unobserved $\mathbf{s}$), and
  \item the normalisation $H_q = \sum_{\mathbf{s}\in\mathcal{O}^{(q)}}h_{\mathbf{s}}^2$
        adapting to the current observation count.
\end{enumerate}

\subsection{Comparison with the fixed-$MN$ convention}

If one instead normalised by $1/(2\pi\sqrt{MN})$ (ignoring masking),
then with $|\mathcal{O}^{(q)}|/MN = p$:
\[
  \mathbb{E}[I^{(q)}_{\mathbf{j}}]\big|_{1/\sqrt{MN}}
  \;\approx\; \frac{p^2\,f(\boldsymbol{\omega}_{\mathbf{j}}) + p(1-p)\sigma^2_Y}{(2\pi)^2}.
\]
The $p^2$ attenuation and $p(1-p)\sigma^2_Y$ white-noise floor are exactly
the spectral distortion terms derived in the companion paper.
The code's $H_q$-normalisation \emph{removes the $p^2$ attenuation}
(since $H_q \approx p\|h\|^2$, dividing by $H_q$ rather than $MN$
compensates by a factor $1/p$), leaving a residual white-noise floor:
\[
  \mathbb{E}[I^{(q)}_{\mathbf{j}}]\big|_{1/\sqrt{H_q}}
  \;\approx\; \frac{f(\boldsymbol{\omega}_{\mathbf{j}}) + [(1-p)/p]\,\sigma^2_Y}{(2\pi)^2}.
\]

\begin{remark}[Practical implication]
The $H_q$ normalisation partially mitigates the $\sigma^2$ downward bias
identified in the companion paper, but does not eliminate the
$(1-p)/p\cdot\sigma^2_Y$ white-noise floor.
For GEMS TCO with $p\approx 0.70$, the residual floor is
$\approx 0.43\,\sigma^2_Y/(2\pi)^2$ per frequency bin.
\end{remark}

\subsection{i.i.d.\ Bernoulli mask: theoretical expectation}

For theoretical analysis, let $m(\mathbf{s})\sim\mathrm{i.i.d.\ Bernoulli}(p)$.
Using $\mathbb{E}[m_\mathbf{u} m_\mathbf{v}] = p^2 + p(1-p)\delta_{\mathbf{u}=\mathbf{v}}$
and the approximation $H_q \approx p\|h\|^2$:

\begin{proposition}[Expected $\tilde{C}_g$ under i.i.d.\ Bernoulli, $H_q$ normalisation]
\label{prop:mask_bernoulli}
\[
\boxed{
  \mathbb{E}_m\!\bigl[\tilde{C}_g(\boldsymbol{\omega}_{\mathbf{j}},\boldsymbol{\omega}_{\mathbf{k}})\bigr]
  \;\approx\;
  \frac{1}{(2\pi)^2}\left[
    \tilde{C}_g^{\mathrm{norm}}(\boldsymbol{\omega}_{\mathbf{j}},\boldsymbol{\omega}_{\mathbf{k}})
  \;+\; \frac{(1-p)\,C(\mathbf{0})}{p\,\|h\|^2}\,F_{h^2}[\mathbf{j}-\mathbf{k}]
  \right],
}
\]
where $\tilde{C}_g^{\mathrm{norm}} = \frac{1}{\|h\|^2}\sum_{\mathbf{u},\mathbf{v}}h(\mathbf{u})h(\mathbf{v})C(\mathbf{u}-\mathbf{v})e^{-i(\omega_j \cdot \mathbf{u} - \omega_k \cdot \mathbf{v})}$
and $F_{h^2}[\mathbf{j}-\mathbf{k}] = \sum_{\mathbf{s}}h_{\mathbf{s}}^2 e^{-i(\boldsymbol{\omega}_{\mathbf{j}}-\boldsymbol{\omega}_{\mathbf{k}})\cdot\mathbf{s}\Delta}$.
\end{proposition}

\begin{remark}[Connection to spectral distortion theory]
At $\mathbf{j}=\mathbf{k}$ and using $F_{h^2}[\mathbf{0}]=\|h\|^2$:
\[
  \mathbb{E}_m\!\bigl[\tilde{C}_g(\omega_j,\omega_j)\bigr]
  \;\approx\;
  \frac{1}{(2\pi)^2}\!\left[f(\omega_j) + \frac{(1-p)}{p}\sigma^2_Y\right].
\]
This is the expected periodogram under $H_q$ normalisation.
Compare with fixed-$MN$ normalisation: the $p^2$ factor is absorbed by
$H_q$, leaving only the $(1-p)/p$ residual floor.
\end{remark}


%───────────────────────────────────────────────────────────────
\section{Three Gaussian Baseline Estimators}
\label{sec:baselines}
%───────────────────────────────────────────────────────────────

All three baselines use Corollary~\ref{cor:pos_freq}:
$\mathrm{Cov}(I_{\mathbf{j}},I_{\mathbf{k}}) \approx 2|\tilde{C}_g(\omega_{\mathbf{j}},\omega_{\mathbf{k}})|^2$,
differing in how $\tilde{C}_g$ is estimated.

\subsection{Baseline 1: Smooth approximation (Priestley, 1981)}

Assume $f(\boldsymbol{\omega})$ is locally flat near $\omega_j,\omega_k$.
Using $f(\omega_m)\approx\sqrt{\hat{f}_j\hat{f}_k}$ in
Proposition~\ref{prop:spectral_conv} and the Hamming identity
$(W^{(q)}\star\overline{W^{(q)}})[\mathbf{d}]\approx H_q\cdot F_{g^2}[\mathbf{d}]$
where $F_{g^2}[\mathbf{d}]=\sum_{\mathbf{s}}g_{\mathbf{s}}^2\,e^{-i(\omega_j-\omega_k)\cdot\mathbf{s}\Delta}$:
\[
  \tilde{C}_g^{\mathrm{smooth}}(\omega_j,\omega_k)
  \;\approx\; \frac{\sqrt{\hat{f}_j\,\hat{f}_k}}{(2\pi)^2}\cdot
    \frac{F_{g^2}[\mathbf{j}-\mathbf{k}]}{H_q},
\]
\[
\boxed{
  \widehat{\mathrm{Cov}}_{\mathrm{smooth}}(I_{\mathbf{j}},I_{\mathbf{k}})
  \;=\; \frac{2\,\hat{f}_{\mathbf{j}}\,\hat{f}_{\mathbf{k}}}{(2\pi)^4}\cdot
    \frac{|F_{g^2}[\mathbf{j}-\mathbf{k}]|^2}{H_q^2}.
}
\]

\begin{remark}
\textbf{Low variance, biased} when $f$ varies across frequencies.
Recommended for trispectrum / non-Gaussian coupling detection.
\end{remark}

\subsection{Baseline 2: Exact full convolution (plug-in)}

Compute $\tilde{C}_g$ directly via Definition~\ref{def:Ctilde}
using empirical $\hat{C}(\mathbf{h})=N_{\mathrm{snap}}^{-1}\sum_n Y_n(\mathbf{s})Y_n(\mathbf{s}+\mathbf{h})$:
\[
  \hat{\tilde{C}}_g(\omega_j,\omega_k)
  \;=\; \frac{1}{(2\pi)^2 H_q}\sum_{\mathbf{u},\mathbf{v}}
    g_{\mathbf{u}}\,g_{\mathbf{v}}\,\hat{C}(\mathbf{u}-\mathbf{v})\,
    e^{-i(\omega_j\cdot\mathbf{u}-\omega_k\cdot\mathbf{v})\Delta},
\]
\[
\boxed{
  \widehat{\mathrm{Cov}}_{\mathrm{exact}}(I_{\mathbf{j}},I_{\mathbf{k}})
  \;=\; 2\,\bigl|\hat{\tilde{C}}_g(\omega_j,\omega_k)\bigr|^2.
}
\]

\begin{remark}[Plug-in bias and variance]
With true $C$: exactly unbiased.
With plug-in $\hat{C}$: plug-in bias present; variance inflated.
Recommended for Whittle bias correction ($\alpha=0$).
\end{remark}

\subsection{Baseline 3: Hybrid (structured shrinkage)}

\[
\boxed{
  \hat{\tilde{C}}_g^{\mathrm{hyb}}(\omega_j,\omega_k)
  \;=\; \alpha\,\tilde{C}_g^{\mathrm{smooth}}(\omega_j,\omega_k)
       + (1-\alpha)\,\hat{\tilde{C}}_g^{\mathrm{exact}}(\omega_j,\omega_k),
}
\]
with covariance estimate $2|\hat{\tilde{C}}_h^{\mathrm{hyb}}|^2$.

\begin{itemize}
  \item $\alpha \to 1$: smooth approx (low variance, biased for non-flat $f$)
  \item $\alpha \to 0$: full convolution (lower bias, high variance)
  \item $\alpha = 0.7$: recommended for trispectrum/physical coupling detection
  \item $\alpha = 0.0$: recommended for Whittle bias correction
\end{itemize}

\subsection{Frequency-pair structure (Hann taper)}

For the outer-product Hann taper, $W[\mathbf{m}]$ has support only at
$\mathbf{m} \in \{0,\pm1\}^2$ per axis. Therefore:
\begin{itemize}
  \item \textbf{Adjacent in one axis only} ($|\mathbf{j}-\mathbf{k}|=1$ in one direction):
        $|\tilde{C}_g(\omega_j,\omega_k)| \approx 0.45\,\sqrt{\hat{f}_j\hat{f}_k}$,
        giving $|\mathrm{Corr}(I_j,I_k)| \approx 0.45$.
  \item \textbf{Diagonal step} (both axes differ):
        $|\tilde{C}_g| \approx 0.01\,\sqrt{\hat{f}_j\hat{f}_k}$,
        trivially near zero.
  \item \textbf{Well-separated}: $\tilde{C}_g \approx 0$.
\end{itemize}
For non-trivial Gaussian baseline validation, pairs must be
\emph{adjacent in exactly one axis}.


%───────────────────────────────────────────────────────────────
\section{No-Taper Case (Rectangular Window)}
\label{sec:no_taper}
%───────────────────────────────────────────────────────────────

When $h \equiv 1$, the tapered cross-covariance reduces to
\[
  \tilde{C}_1(\omega_j,\omega_k)
  = \frac{1}{MN}\sum_{\mathbf{u},\mathbf{v}} C(\mathbf{u}-\mathbf{v})\,
    e^{-i(\omega_j\cdot\mathbf{u}-\omega_k\cdot\mathbf{v})}
  = f(\omega_j)\,\delta_{\mathbf{j}=\mathbf{k}},
\]
since at distinct Fourier frequencies the sum cancels exactly.
Therefore, for $\mathbf{j} \neq \mathbf{k}$ (and $\mathbf{j} \neq -\mathbf{k}$):
\[
\boxed{
  \mathrm{Cov}(I_{\mathbf{j}},\, I_{\mathbf{k}}) = 0
  \quad\textbf{(exact, no approximation, no taper)}
}
\]

\begin{remark}
The sample mean correction $g(\mathbf{s}) = Y(\mathbf{s}) - \bar{Y}$ introduces
a minor perturbation at the DC frequency, negligible for large $M,N$.
\end{remark}


%───────────────────────────────────────────────────────────────
\section{Non-Gaussian Residual (Trispectrum)}
\label{sec:nongauss}
%───────────────────────────────────────────────────────────────

\subsection{Empirical covariance}

Given $N_{\mathrm{snap}}$ snapshots with periodogram values $I_n(\boldsymbol{\omega}_{\mathbf{j}})$:
\[
  \widehat{\mathrm{Cov}}_{\mathrm{emp}}(I_{\mathbf{j}},I_{\mathbf{k}})
  \;=\; \frac{1}{N_{\mathrm{snap}}}\sum_{n=1}^{N_{\mathrm{snap}}}
    \bigl(I_n(\boldsymbol{\omega}_{\mathbf{j}}) - \hat{f}_{\mathbf{j}}\bigr)
    \bigl(I_n(\boldsymbol{\omega}_{\mathbf{k}}) - \hat{f}_{\mathbf{k}}\bigr).
\]

\subsection{Trispectrum residual}

\[
  \hat{\kappa}_{\mathbf{j},\mathbf{k}}
  \;=\; \widehat{\mathrm{Cov}}_{\mathrm{emp}} - 2\,|\hat{\tilde{C}}_h(\omega_j,\omega_k)|^2.
\]

\subsection{Full decomposition}

\[
  \underbrace{\widehat{\mathrm{Cov}}_{\mathrm{emp}}}_{\text{total}}
  \;=\;
  \underbrace{2\,|\hat{\tilde{C}}_h(\omega_j,\omega_k)|^2}_{\text{Gaussian (Wick)}}
  \;+\;
  \underbrace{\hat{\kappa}_{\mathbf{j},\mathbf{k}}}_{\text{non-Gaussian (trispectrum)}}
\]


%───────────────────────────────────────────────────────────────
\section{Significance Threshold}
\label{sec:threshold}
%───────────────────────────────────────────────────────────────

Under the Whittle independence assumption, sample correlation
$\hat\rho_{\mathbf{jk}}$ is asymptotically $\mathcal{N}(0, 1/N_{\mathrm{snap}})$.
Two-sided 5\% threshold:
\[
  |\hat\rho_{\mathbf{jk}}| > \frac{2}{\sqrt{N_{\mathrm{snap}}}}.
\]
For $N_{\mathrm{snap}} = 248$: threshold $\approx 0.127$.


%───────────────────────────────────────────────────────────────
\section{Gaussian Simulation Validation}
\label{sec:validation}
%───────────────────────────────────────────────────────────────

The formula
$\mathrm{Cov}(I_j, I_k) = 2|\tilde{C}_g(\omega_j,\omega_k)|^2$
was validated by two simulations.

\subsection{1-D AR(1) validation}

Setting: $N=128$, $T=2000$ realizations, $\rho_{\mathrm{AR}}=0.85$,
Hann taper. True spectrum $f_{\mathrm{true}}[\omega_j]
= \sigma^2/|1-\rho_{\mathrm{AR}}e^{-i\omega_j}|^2$.
For adjacent pairs ($|\mathbf{j}-\mathbf{k}|=1$), the formula predicts
$|\mathrm{Corr}| \approx 0.44$--$0.45$; empirical match within
$2/\sqrt{T} \approx 0.045$.
\textbf{Formula confirmed.}

\subsection{2-D spatial field validation}

Setting: $M\times N = 114\times 159$, $N_{\mathrm{snap}}=500$, Hann taper,
exponential covariance. Gaussian field via circulant embedding:
\[
  Y = \sqrt{2}\cdot\mathrm{Re}\!\left(\mathrm{IFFT}\!\left(\sqrt{f_{\mathrm{true}}}\cdot Z\right)\right)\cdot\sqrt{MN},
  \quad Z\sim\mathcal{CN}(0,I),
\]
where $\sqrt{2}$ corrects the variance halving from $\mathrm{Re}(\cdot)$.

Frequency groups:
\begin{itemize}
  \item \textbf{NearLon}: $(5,8),(5,9),(5,10)$ --- same lat-index, adjacent lon-index.
        Baseline predicts $|\tilde{C}_g| \approx 0.45\sqrt{f_jf_k}$,
        $|\mathrm{Corr}| \approx 0.45$.
  \item \textbf{NearLat}: $(8,20),(9,20),(10,20)$ --- adjacent lat-index,
        same lon-index. Same prediction.
  \item \textbf{Far}: $(20,40),(35,60),(50,75)$ ---
        $|\tilde{C}_g| \approx 0$, $|\mathrm{Corr}| \approx 0$.
\end{itemize}
Empirical results match formula for both Near groups.
\textbf{Formula confirmed in 2-D.}

\begin{remark}[Frequency selection pitfall]
``Diagonal step'' pairs ($\mathbf{j}-\mathbf{k}$ differs in both axes) give
$|\tilde{C}_g| \approx 0.01\sqrt{f_jf_k}$ trivially, making off-diagonal
validation vacuous.
Pairs adjacent in exactly one axis are required for non-trivial
validation of $\tilde{C}_g$.
\end{remark}


%───────────────────────────────────────────────────────────────
\section{Summary}
\label{sec:summary}
%───────────────────────────────────────────────────────────────

\subsection{Key formula hierarchy}

\begin{center}
\begin{tabular}{ll}
\toprule
Object & Expression \\
\midrule
Effective weight & $g^{(q)}_\mathbf{s} = h_\mathbf{s}\cdot\mathbf{1}[\mathbf{s}\in\mathcal{O}^{(q)}]$,\quad $H_q=\sum_\mathbf{s}g_\mathbf{s}^2$ \\[4pt]
Tapered DFT & $J^{(q)}(\omega_k) = \frac{1}{2\pi\sqrt{H_q}}\mathcal{F}_{2D}[g^{(q)}\odot X^{(q)}](\mathbf{k})$ \\[4pt]
Cross-cov kernel & $\tilde{C}_g(\omega_j,\omega_k) = \frac{1}{(2\pi)^2 H_q}\sum_{u,v}g_u g_v C(u-v)e^{-i(\omega_j u - \omega_k v)\Delta}$ \\[4pt]
Gauss cov & $\mathrm{Cov}(I_j,I_k) = 2|\tilde{C}_g(\omega_j,\omega_k)|^2$ \\[4pt]
$H_q$ mask effect ($j=k$) & $\mathbb{E}[\tilde{C}_g(\omega_j,\omega_j)] \approx \frac{1}{(2\pi)^2}[f(\omega_j) + \frac{1-p}{p}\sigma^2_Y]$ \\
\bottomrule
\end{tabular}
\end{center}

\subsection{Three Gaussian baseline estimators}

\begin{table}[h!]
\centering
\begin{tabular}{lcccc}
\toprule
Method & $\tilde{C}_g$ estimate & Bias & Variance & Recommended use \\
\midrule
Smooth & $\sqrt{\hat{f}_j\hat{f}_k}\cdot F_{h^2}[j-k]/(MN)$
  & Biased & Low & Non-Gaussian detection \\
Exact & $\frac{1}{MN}\sum_{u,v}h(u)h(v)\hat{C}(u-v)e^{-i(\cdot)}$
  & Plug-in & High & Whittle bias correction \\
Hybrid & $\alpha$-mixture & Intermediate & Intermediate & Physical coupling \\
No taper & $f(\omega_j)\delta_{j=k}$ & None & Low & Baseline comparison \\
\bottomrule
\end{tabular}
\end{table}


%───────────────────────────────────────────────────────────────
\section{Empirical Findings (GEMS TCO, July 2024)}
\label{sec:empirical}
%───────────────────────────────────────────────────────────────

Grid: $114\times 159$, $D=31$ days, $N_{\mathrm{snap}}=248$ snapshots.
Observation fraction $p \approx 0.70$--$0.75$ (missing $\approx 25$--$30\%$).
Significance threshold $2/\sqrt{248}\approx0.127$.

\paragraph{Masked kernel.}
Under the actual GEMS masking, $\tilde{C}_{h,m}$ satisfies
(Proposition~\ref{prop:mask_bernoulli}):
\[
  \mathbb{E}_m[\tilde{C}_{h,m}(\omega_j,\omega_k)]
  = p^2\,\tilde{C}_g(\omega_j,\omega_k)
  + \frac{p(1-p)\sigma^2_Y}{MN}\,F_{h^2}[j-k].
\]
At $j=k$: $\approx 0.49 f(\omega_j) + 0.21\sigma^2_Y$ (for $p=0.70$).
The second term inflates the effective diagonal cross-covariance and
must be accounted for when using the exact baseline.

\paragraph{Frequency bands and findings.}

\begin{table}[h!]
\centering
\begin{tabular}{lll}
\toprule
Band & Gaussian baseline $|\tilde{C}_g|$ & Empirical finding \\
\midrule
Near (adj.\ in one axis) & $\approx 0.45\sqrt{\hat{f}_j\hat{f}_k}$
  & Excess above 0.45 $\Rightarrow$ non-Gaussian \\
Near $\times$ Mid/High & $\approx 0$
  & Any $>0.127$ $\Rightarrow$ non-Gaussian coupling \\
Mid $\times$ Mid & $\approx 0$
  & Previously: Gaussian sufficient \\
High $\times$ High & $\approx 0$
  & $|\hat\rho|\leq0.18$; moderate non-Gaussian \\
\bottomrule
\end{tabular}
\end{table}

\paragraph{Main conclusion.}
The Whittle independence assumption is most reliable for Mid-frequency
bands ($\lambda\approx 0.3$--$0.6^\circ$).
High-frequency bands show moderate non-Gaussian coupling
($|\hat\rho|\approx0.18$).
Zero-imputation masking inflates the effective diagonal of
$\tilde{C}_{h,m}$ by $p(1-p)\sigma^2_Y$, which should be corrected
before estimating the Gaussian baseline.

\bigskip
\begin{thebibliography}{9}
\bibitem{Brillinger1975}
Brillinger, D.R.\ (1975).
\emph{Time Series: Data Analysis and Theory}.
Holt, Rinehart and Winston.

\bibitem{Priestley1981}
Priestley, M.B.\ (1981).
\emph{Spectral Analysis and Time Series}.
Academic Press.

\bibitem{Sykulski2019}
Sykulski, A.M., Olhede, S.C., Guillaumin, A.P., Lilly, J.M., and Early, J.J.\
(2019).
The debiased Whittle likelihood.
\emph{Biometrika}, \textbf{106}(2), 251--266.
\end{thebibliography}

\end{document}
